\section{OBJECTIVE AND REVIEW OF ALL VERSIONS}
\textbf{Erd\H{o}s \#789; independent attempt, 2026-09-23.}
Call a finite integer set admissible if equal sums of nonempty subsets have equal cardinalities. Determine the asymptotic order of the largest $h(n)$ guaranteed as an admissible subset of every $n$-element integer set. I read both \texttt{789.tex} and \texttt{789\_v2.tex} completely, including the first file's unrelated archived discussion of \#873 and \#875. The relevant arguments transfer an interval upper bound to $h(n)$ but do not close the known lower/upper exponent gap. The checked formal statement confirms that sums are subset sums, without repeated summands.
\section{WORK: TRANSLATION CAN REMOVE EVERY OBSTRUCTION}
Let $B=\{b_1,\ldots,b_m\}\subset\mathbb Z$, $L=\min B$, and $D=\max B-L$. If an integer $T$ satisfies
\[
L+T>(m-1)D,
\]
then $B+T$ is admissible. Indeed put $u=L+T>0$. A sum of $r$ elements lies in $[ru,r(u+D)]$. For $0\le r<m$,
\[
r(u+D)<(r+1)u
\]
because $rD\le(m-1)D<u$. Thus the possible sum intervals for different cardinalities are disjoint. This even proves the stronger convention allowing the empty subset.
There are in fact only finitely many exceptional translations, in either direction. Cancel the common elements of two unequal-cardinality subsets. A remaining collision in $B+T$ has the form
\[
\sum_{i=1}^m\varepsilon_i b_i+T\sum_{i=1}^m\varepsilon_i=0,
\qquad \varepsilon_i\in\{-1,0,1\},\quad\sum_i\varepsilon_i\ne0.
\]
Each sign vector permits at most one integer $T$. Vectors $\varepsilon$ and $-\varepsilon$ give the same candidate, so there are at most $(3^m-1)/2$ exceptional integer translations. This is only an upper bound: many candidates are nonintegral or coincide.
\section{ADVERSARIAL VERIFICATION AND REMAINING GAP}
Admissibility is therefore not translation invariant. For example $\{1,2,3\}$ is not admissible since $1+2=3$, whereas $\{11,12,13\}$ is admissible by the proved interval criterion. This prevents normalizing an arbitrary input set by translation and then carrying an admissible subset back unchanged. Nonzero scaling does preserve the property; translation adds a cardinality-dependent term and does not.
The universal problem concerns the input at its actual location, so the translation lemma cannot improve $h(n)$. Using $[1,n]$ as a particular adverse input, as the supplied upper-bound argument does, remains valid. The useful conclusion here is a precise obstruction to a plausible but invalid compression strategy, not a new asymptotic estimate.
\section{SOURCES CHECKED}
Both complete local versions; official indexed entry \url{https://www.erdosproblems.com/789}; the subset-based definition in \url{https://raw.githubusercontent.com/google-deepmind/formal-conjectures/main/FormalConjectures/ErdosProblems/789.lean}.
\section{FINAL}
\textbf{UNRESOLVED.} The effect of translation is proved exactly; the asymptotic order of $h(n)$ remains undetermined.
\section{COMPLETION ESTIMATE}
Conservative subjective progress, not a probability of eventual success.
\textbf{COMPLETION: 8\%}
